\section{Enfoque de desarrollo seleccionado y justificación}

Para un proyecto de solo dos semanas, no es práctico planificar todo al inicio como en el modelo tradicional en cascada. Se eligió un enfoque \textbf{adaptativo e incremental}, combinando prácticas sencillas de Scrum, Kanban y XP:

\subsection{Prácticas adoptadas del proceso}

\begin{itemize}[nosep]
    \item \textbf{Aportes de Scrum:} Trabajo dividido en ciclos cortos (sprints) con una meta clara de entrega. Se realizan reuniones breves de 10 minutos para revisar el avance diario y una reunión al final del ciclo para ver qué mejorar.
    \item \textbf{Aportes de Kanban:} Uso de un tablero en Jira con columnas visibles y un límite de trabajo en curso (WIP). Esto ayuda a que el equipo no abra muchas tareas a la vez, aplicando la regla de \textit{terminar antes de empezar algo nuevo}.
    \item \textbf{Aportes de XP:} Pruebas automáticas con PHPUnit en las funciones más delicadas (control de cupos y DNI) y revisión obligatoria de código entre compañeros en GitHub antes de unir cambios a la rama principal.
    \item \textbf{Uso de IA en el proceso:} La Inteligencia Artificial se utiliza como asistente para redactar historias de usuario o consultar sintaxis en Laravel. Cada sugerencia de la IA es revisada manualmente y comprobada con pruebas locales antes de usarse.
\end{itemize}
